The idea of discussion is to bring up points that you noticed during research.
As a chapter, this is the most freeform one.
A good way to discuss the points is first to write up subtitles with clear statements and then expand on those.
Discussion points can occasionally be also in a question format or refer to a larger context beyond the thesis to bring in ideas from other papers to consider.
In case your work has limitations that you can observe, it can be a good idea to write a separate section considering those.

% =========================================================================
% DRAFT NOTES: Prompt Caching Latency Inversion in Enterprise RAG (Nordea Closed)
% =========================================================================
%
% -------------------------------------------------------------------------
% 1. EMPIIRINEN HAVAINTO (Latenssi-inversio Enterprise Closed -aineistossa, N=100)
% -------------------------------------------------------------------------
% - Vastoin intuitiota ja perusoletusta, Module M3 (Prompt Caching) ei tuo 
%   latenssihyötyjä yrityskontekstissa, vaan nostaa kokonaislatenssia:
%     * C0 Baseline:    T_total P50 = 1 407 ms | TTFT P50 = 1 135 ms
%     * C3 (M3 Prompt Caching): T_total P50 = 1 714 ms (+21.8 %) | TTFT P50 = 1 563 ms (+37.7 %)
%     * C5 (M1+M2 S2B + SR):    T_total P50 = 1 446 ms | TTFT P50 = 1 220 ms
% - C3 on merkittävästi hitaampi kuin C0 Baseline ja jopa selvästi hitaampi kuin 
%   C5 (S2B + SR), jonka odotettiin olevan raskas ylimääräisen semanttisen rerankkauksen 
%   vuoksi (rerank-vaiheen P50 = 108 ms).
%
% -------------------------------------------------------------------------
% 2. TEKNINEN JUURISYY: T_prefill ja etä-KV-cachen overhead
% -------------------------------------------------------------------------
% - Kun puretaan latenssikomponentit (context_inspector.py ja evaluoidut ajot):
%     * T_prefill C0: median = 827 ms
%     * T_prefill C3: median = 1 308 ms (+481 ms)
%     * T_stream C0:  median = 147 ms
%     * T_stream C3:  median = 105 ms
% - Itse vastaustokenien generointi (T_stream) on C3:ssa identtinen tai jopa 
%   aavistuksen nopeampi kuin C0:ssa. Koko latenssirangaistus syntyy mallin 
%   prefill-vaiheessa ennen ensimmäisen tokenin tuloa (TTFT).
% - Gemini 3.5 Flash laskee Googlen TPU-laitteistolla ~7 000 - 8 500 tokenin rinnakkaisen 
%   prefill-huomiolaskennan (attention prefill) muutamissa kymmenissä millisekunneissa.
% - Kun pyyntö käyttää Vertex AI explicit context cachingia (cachedContentName):
%     1. Vertex AI -infrastruktuuri joutuu validoimaan välimuistiviitteen ja etsimään 
%        sen hajautetusta KV-välimuistivarastosta.
%     2. Pyyntö on reititettävä tai sarjallistettu KV-tila on ladattava verkon yli 
%        juuri sille TPU-kiihdyttimelle, joka pyyntöä palvelee.
%     3. Tämä etätilan sidonta ja välimuistinhaku aiheuttaa kiinteän ~300-500 ms 
%        RPC- ja synkronointiviiveen (fixed coordination overhead).
%
% -------------------------------------------------------------------------
% 3. KONTEKSTIKOON "BREAK-EVEN" -KYNNYS
% -------------------------------------------------------------------------
% - Välimuistiin tallennettu staattinen etuliite (järjestelmäohjeet + few-shot -kirjasto, 
%   prompts/answer-generation.st) on tarkalleen 7 157 tokenia.
% - Vertex AI / Gemini Context Caching on suunniteltu massiivisille konteksteille 
%   (32k - 1M+ tokenia, esim. kirjat, koodirepot, videot). Tällöin KV-uudelleenlaskenta 
%   kestäisi sekunteja, jolloin ~300-500 ms noutoviive tuo suuren nettovoiton.
% - 7 000 tokenin kokoluokassa etävälimuistin noutaminen kestää pidempään kuin se aika, 
%   joka kuluisi tokenien puhtaaseen uudelleenlaskentaan TPU:n HBM-muistissa tyhjästä.
%   -> Yritys-RAG:ssa syntyy "latency break-even" -ongelma: konteksti on liian lyhyt 
%      latenssisäästöille, mutta ylittää palveluntarjoajan minimivälimuistikynnyksen (4 096 tok).
%
% -------------------------------------------------------------------------
% 4. KONTEKSTIVÄLIMUISTI VS. SEMANTTINEN VASTAUSVÄLIMUISTI
% -------------------------------------------------------------------------
% - M3 toteuttaa mallitason KV-kontekstivälimuistin (Gemini CachedContent), 
%   EI sovellustason semanttista vastausvälimuistia (kuten GPTCache tai Redis).
% - Vastausvälimuisti toisi osumalla 0-10 ms latenssin, mutta kontekstivälimuistissa 
%   LLM joutuu aina suorittamaan uuden päättelyn ja generoinnin dynaamisesti 
%   noudetuille yritysdokumenteille.
%
% -------------------------------------------------------------------------
% 5. DIPLOMITYÖN TULOS & PARETO-ANALYYSI (RQ1 vs. RQ2)
% -------------------------------------------------------------------------
% - Prompt Caching ei toimi yritystason RAG:ssa latenssinoptimoijana, vaan puhtaana 
%   kustannusoptimointina:
%     * Latenssi (RQ1): T_total nousee +21.8 % (1 407 ms -> 1 714 ms).
%     * Kustannus (RQ2): Ajokustannus putoaa -70.7 % ($13.58 -> $3.98 / 1k kyselyä; 
%       yhdessä M4-reitityksen kanssa jopa $2.54 / 1k eli -81.3 %).
% - Laatu (RQ3): F1 säilyy identtisenä (C0: 59.4 % vs. C3: 58.8 %).
% - Johtopäätös: Järjestelmäarkkitehdin on tehtävä tietoinen Pareto-valinta: 
%   ollaanko valmiita maksamaan ~300 ms lisälatenssia käyttäjälle vastineena yli 
%   70 % operatiivisesta pilvikustannussäästöstä.
%
% =========================================================================
% DRAFT NOTES 1B: Prompt Caching Häntälatenssin Sulakkeena (Circuit-Breaker)
% =========================================================================
% - Vaikka PC nostaa mediaanilatenssia, MuSiQuessa se toimii kriittisenä 
%   P95-häntälatenssin leikkurina monivaiheisissa ketjuissa:
%     * S2B + SR (C5):    P95 = 7 771 ms -> S2B + SR + PC (C11):    P95 = 2 202 ms (-71.7 %)
%     * S2B + MC (C7):    P95 = 14 154 ms -> S2B + PC + MC (C13):   P95 = 4 644 ms (-67.2 %)
%     * S2B + SR + MC (C12): P95 = 30 883 ms -> S2B + SR + PC + MC (C15): P95 = 4 306 ms (-86.1 %)
%     * S2B + PC (C6):    P50 = 1 250 ms (-8.0 %, p < 0.001) ja P95 = 1 745 ms (baseline: 3 397 ms)
% - JSONL-raakadata paljastaa: Ilman välimuistia monimutkaiset ketjut kärsivät satunnaisista 
%   20-90 s TPU-muistinvaraus- ja jonotuspiikeistä ruuhka-aikoina (esim. C12 kyselyssä 
%   3hop1__135793_87694_64412: T_prefill = 89.7 s).
% - Cachen käyttö eliminoi nämä ääripiikit, koska 7 157 tokenin esilaskettu tila on jo 
%   valmiiksi kiinnitetty TPU:n muistiin. PC toimii siis SLA-sulakkeena (Circuit-Breaker).
%
% =========================================================================
% DRAFT NOTES 2: Frontier-mallin (Gemini 3.1 Pro) Over-Refusal Paradox (RQ3 / M4)
% =========================================================================
% - Reititysjakautuma:
%     * MuSiQue: Pro-dominantti (77 % Pro, 21 % Flash, 2 % Lite; 3/4-hop 100 % Pro).
%     * Enterprise Closed: Lite/Flash-dominantti (49 % Lite, 38 % Flash, 13 % Pro; 
%       1-hop 97.5 % Lite, 2-hop 67.5 % Flash / 30 % Pro).
% - Paradoksi MuSiQuessa: Frontier-malli Gemini 3.1 Pro kieltäytyy vastaamasta 
%   vastattaviin kysymyksiin huomattavasti herkemmin kuin Gemini 3.5 Flash:
%     * C0 Baseline (Flash): FRR = 42.5 %
%     * C4 MC (Pro-dominantti): FRR = 63.9 % (Pron oma FRR vastattavissa = 63.9 %)
%     * C14 SR+PC+MC: FRR = 78.3 % Prolla (koko putken FRR = 68.8 %)
% - Juurisyy JSONL-datassa (Ehdollinen Context Recall -analyysi):
%     * Kun haettu konteksti on osittainen (CR < 1.0, eli 1 välihyppy puuttuu top-5:stä):
%       Pro kieltäytyy 87.0 % - 96.8 % ajasta, kun taas Flash kieltäytyy vain 50.0 % - 68.6 % 
%       ja kykenee päättelemään puuttuvan askeleen.
%     * Täydelläkin kontekstilla (CR = 1.0) Pro kieltäytyy 42.3 % - 58.6 % ajasta monimutkaisissa 
%       monihyppytehtävissä tulkiten synteesin epävarmaksi (Flash kieltäytyy vain 0 - 22 %).
% - Johtopäätös: Pron tiukka RLHF- ja grounding-ohjeistus tekee siitä tuotantoon turvallisen 
%   (100 % OOD Refusal Accuracy), mutta se rankaisee monihyppyhaun osittaista osumatarkkuutta.
%
% =========================================================================
% DRAFT NOTES 3: Synergia-bifurkaatio: SR + MC (C9) romahdus vs. S2B + MC (C7) menestys
% =========================================================================
% - Enterprise Closedissa havaitaan jyrkkä laatuero:
%     * S2B + MC (C7):  Q = 0.6587 (toiseksi paras), Faith = 0.858, FRR = 11.2 %
%     * SR + MC (C9):   Q = 0.5418 (koko testin huonoin, -12.7 %), Faith = 0.780, FRR = 17.5 %
%     * Huom: C9:llä on silti koko testin KORKEIN Context Precision (0.683 vs. C7:n 0.661)!
% - Juurisyy JSONL-datasta (Kyselyt 1hop__nordea_0030, 2hop__nordea_0009, 0019, 0020, 0040):
%     * Re-ranker ilman S2B-laajennusta karsii pois reunaehtojen ja 2-hypyn taustakappaleet 
%       (kontekstin pituus putoaa 4.5k merkkiin vs. C7:n 8.1k merkkiin).
%     * Flash-Lite ei kykene rekonstruoimaan puuttuvia ehtoja kapeista paloista: se kieltäytyy 
%       tai hallusinoi ehtoja (Faithfulness putoaa Litellä 0.906 -> 0.849, esim. kysely 0040: Faith = 0.50).
%     * Pro kieltäytyy herkästi puutteellisella kontekstilla (Pron FRR C9:ssä peräti 50.0 %!).
%     * C7:ssä S2B:n 1 024 tokenin vanhempidokumentti toimii suojakilpenä: Litellä FRR on 
%       vain 2.4 % ja Faithfulness 0.906. S2B suojaa kevyitä ja tiukkoja malleja pirstaloitumiselta.
%
% =========================================================================
% DRAFT NOTES 4: Hakukohinan Kynnysarvo (Retrieval Noise Threshold) & Korpusten Inversio
% =========================================================================
% - Suljetussa korpuksessa (Enterprise Closed, 20 kandidaattia / kysely, 1 701 palaa):
%     * C1 S2B pärjää yksin erinomaisesti (Q = 0.6449 ns).
%     * Re-rankingin lisääminen (C5 S2B + SR) LASKEE laatua (Q = 0.5813 ns), koska puhdas 
%       kanta ei tarvitse karsintaa ja re-ranker rikkoo peräkkäisten lapsipalojen koheesion.
% - Avoimessa raakakorpuksessa (Enterprise Raw, 2 420 uniikkia palaa, PDF-ehdot ja web):
%     * Baseline C0 romahtaa (Q = 0.5022, 2-hop FRR hyppää 12.5 % -> 27.5 %).
%     * Re-ranking toimii välttämättömänä kohinasuodattimena (Noise Gate), joka siivoaa distraktorit 
%       ennen S2B-laajennusta:
%         - S2B + SR (C5):    Q = 0.6433** (p_adj = 0.0077, +28.1 % parannus baselineen)
%         - S2B + SR + PC (C11): Q = 0.6466** (p_adj = 0.0029, +28.8 % parannus)
%         - S2B + SR + MC (C12): Q = 0.6162*  (p_adj = 0.0470, +22.7 % parannus)
%         - Context Precision paranee tilastollisesti merkitsevästi (p_adj < 0.05) peräti 7 S2B-putkessa.
%         - 2-hop FRR putoaa 27.5 %:sta 7.5 %:iin!
% - Johtopäätös: Re-ranking tuo tilastollisesti vahvistetun hyödyn vasta, kun hakukohina 
%   ylittää kriittisen kynnysarvon. Kuratoidussa kannassa se on redundanssi.
%
% =========================================================================
% DRAFT NOTES 5: Harmonisen Keskiarvon Q (eps = 10^-4) "Heikoimman Lenkin Rangaistus"
% =========================================================================
% - Säännöllistetyn harmonisen keskiarvon kaava (Eq. 4.25):
%     * Jos YKSIKIN neljästä metriikasta (CP, CR, Faith, AR) on 0.0:
%       nimittäjään tulee 1/10^-4 = 10 000, jolloin Q(q) romahtaa välittömästi arvoon <= 0.0004.
% - Empiirinen todiste: MuSiQuessa 68.8 % - 77.5 % vastattavista kyselyistä sisältää vähintään 
%   yhden nollametristen (useimmiten FRR:n aiheuttama Faith=0, AR=0 tai noudon CR=0).
% - Tilastollinen seuraus: Koska yli puolet otoksesta on nollassa, mediaani-Q naulautuu arvoon 
%   0.0002 kaikilla putkilla, ja parittainen mediaaniero muuttuu muotoon mu_delta = 0.0000.
% - Tämä nollaa parametrittoman Wilcoxon-testin tilastollisen voiman (sidokset / ties).
% - Metodologinen opetus: Harmoninen Q on ankara teollinen luotettavuusmittari (penalizes single-point-of-failure), 
%   mutta tilastolliseen hypoteesitestaukseen se vaatii suuren otoskoon tai jatkuvamman metriikan.
%
% =========================================================================
% DRAFT NOTES 6: Poikkileikkaavat Havainnot (PC Attention Anchoring & S2B Latenssi)
% =========================================================================
% - PC:n huomion ankkurointivaikutus (Attention Anchoring):
%     * C6 (S2B + PC) saavuttaa korkeamman Faithfulnessin (0.886) kuin C1 S2B (0.863).
%     * Staattisen kehotteen ja CoT-esimerkkien jäädyttäminen KV-cacheen ankkuroi huomion, 
%       estäen tarkkaavaisuuden hajoamista (attention dilution / drift) pitkissä 1 024 tokenin ikkunoissa.
% - Small-to-Bigin negatiivinen latenssidelta yritysaineistoissa (-4.3 % Closed, -2.4 % Raw):
%     * PostgreSQL:n parent-laajennus vie vain 1-2 ms (T_expand).
%     * Kompaktimpi 128 tokenin lapsi-indeksi leikkaa noudon P95-häntälatenssin puoleen (674 -> 317 ms).
%     * Eheä parent-konteksti nopeuttaa LLM:n striimausta (T_stream P50: 147 -> 131.5 ms), koska 
%       malli vastaa varmemmin ilman epävarmuuden kiertelyä.
%

